\documentclass[11pt]{article}
\usepackage[a4paper,margin=30mm]{geometry}
\usepackage{amsmath,amssymb,amsthm,mathtools,array}
\usepackage[hidelinks]{hyperref}
\usepackage{xcolor}

\newtheorem{theorem}{Theorem}
\newtheorem{lemma}[theorem]{Lemma}
\newtheorem{proposition}[theorem]{Proposition}
\newtheorem{corollary}[theorem]{Corollary}
\theoremstyle{definition}
\newtheorem{definition}[theorem]{Definition}
\newtheorem{hypothesis}[theorem]{External interface}
\newtheorem{unknown}[theorem]{UNKNOWN obligation}
\theoremstyle{remark}
\newtheorem{remark}[theorem]{Remark}

\newcommand{\R}{\mathbb R}
\newcommand{\T}{\mathbb R/\mathbb Z}
\newcommand{\mesh}{\operatorname{mesh}}
\newcommand{\Ind}{\operatorname{Ind}}
\newcommand{\diam}{\operatorname{diam}}
\newcommand{\status}[1]{\par\smallskip\noindent\textbf{Status:}
  \begingroup\small\sffamily #1.\endgroup\par\smallskip}

\title{Formal dependency draft for the vanishing critical
  $2$-variation prescribed-angle rectangle claim}
\author{Issue \#134; companion specification for PR \#122}
\date{Draft of 30 August 2026}

\begin{document}
\maketitle

\begin{abstract}
This document is a formalization interface, not an unconditional Square Peg
claim.  It isolates the finite-partition and Liouville-primitive arguments
that can be proved from the displayed definitions, states the external
theorems with exactly the interfaces used, and records every unresolved
dependency.  In particular, no positive-size square and no four-distinct-
vertices conclusion is asserted unless the nondegenerate conclusion of the
Asano--Ike theorem is independently verified from the primary source.
\end{abstract}

\section{Status vocabulary and theorem interface}

Every lemma below carries one of the following statuses.
\begin{itemize}
\item \texttt{PROVED}: proved in this document from earlier displayed
  definitions and stated standard facts.
\item \texttt{CITED}: imported from a named external source; citation
  fidelity remains part of the verification obligation.
\item \texttt{STANDARD-CITED}: a standard theorem is used, but not reproved
  foundationally here.
\item \texttt{STANDARD-ASSUMED}: a precise auxiliary interface is stated,
  but its formal proof and source pinning are outside this draft.
\item \texttt{UNKNOWN}: the exact statement or a required construction has
  not been certified sufficiently to use as an assumable premise.
\end{itemize}

The intended dependency chain is
\[
\begin{gathered}
 V_0\longrightarrow\text{polygonal approximation}
 \longrightarrow\text{regular smoothing}\\
 \longrightarrow\text{uniform local }2\text{-variation}
 \longrightarrow\text{primitive compactness}
 \longrightarrow\text{Asano--Ike conclusion}.
\end{gathered}
\]
The local consequences of $V_0$, the fourth box, and the elementary parts
of the fifth box are proved below.  The other boxes expose imported or
unknown interfaces.

\section{\texorpdfstring{Vanishing $2$-variation}{Vanishing 2-variation}
and polygonal error}

Let $c\colon\R\to\R^2$ be continuous, $1$-periodic, and injective on
$[0,1)$.  For an interval $I=[a,b]$, define
\[
 \|x\|_{2\text{-var};I}^2
 :=\sup_{a=t_0<\cdots<t_m=b}
       \sum_{j=0}^{m-1}|x(t_{j+1})-x(t_j)|^2.
\]
For a partition $D=\{0=t_0<\cdots<t_m=1\}$ put
\[
 \mesh(D)=\max_j(t_{j+1}-t_j),\qquad
w_2(c,\delta)^2
 :=\sup_{\mesh(D)\leq\delta}
       \sum_j|c(t_{j+1})-c(t_j)|^2.
\]
Both suprema initially take values in $[0,\infty]$.

\begin{definition}[Critical vanishing-variation hypothesis]
Write $V_0(c)$ for
\[
 \exists\delta_0>0:\ w_2(c,\delta_0)<\infty
 \quad\text{and}\quad w_2(c,\delta)\longrightarrow0
                 \quad(\delta\downarrow0).
\]
Thus every real-valued variation calculation below occurs in the finite
regime $0<\delta\leq\delta_0$.
\end{definition}

\begin{lemma}[Lifted short-interval modulus]\label{lem:local-modulus}
If $V_0(c)$ holds, $0<\delta\leq\tfrac12$, and $I\subset\R$ has
$|I|\leq\delta$, then
\[
       \|c\|_{2\text{-var};I}\leq\sqrt2\,w_2(c,\delta).
\]
\status{PROVED}
\end{lemma}

\begin{proof}
If $I$ does not cross an integer, extend any partition of $I$ to a
full-period partition of mesh at most $\delta$.  If it crosses one integer,
insert the seam.  The insertion replaces at most one increment $u+v$ by
$u,v$, and
\(
 |u+v|^2\leq2(|u|^2+|v|^2).
\)
After reducing the two pieces modulo one, fill every complementary gap by
subintervals of length at most $\delta$ and combine the nodes into a
full-period partition.  Since $|I|\leq\delta\leq1/2$, at most one integer
seam occurs.  Take suprema.
\end{proof}

For a partition $D$, let $c^D$ be the affine interpolant of $c$ on every
interval of $D$.

\begin{lemma}[Every-partition interpolation estimate]
For every partition $D$ of $[0,1]$,
\[
 \|c-c^D\|_{2\text{-var};[0,1]}
 \leq \sqrt8\,w_2(c,\mesh(D)).                         \tag{2.1}
\]
Consequently $V_0(c)$ implies $c^D\to c$ in $2$-variation along
\emph{every} sequence with $\mesh(D)\to0$.
\status{PROVED; replaces the load-bearing Yang quantifier}
\end{lemma}

\begin{proof}
Write $I_j=[t_j,t_{j+1}]$ and $y=c-c^D$.  Since $y$ vanishes at every
$t_j$, inserting all partition nodes into an arbitrary variation partition
costs at most a factor $2$ in the squared sum.  Indeed, although one
crossing increment may meet many nodes of $D$, $y=0$ at every inserted
node, so only its two endpoint pieces can be nonzero.  Hence
\[
 \|y\|_{2\text{-var}}^2
 \leq2\sum_j\|y\|_{2\text{-var};I_j}^2.
\]
On $I_j$, Minkowski's inequality and the fact that the affine chord has
$2$-variation equal to $|c(t_{j+1})-c(t_j)|$ give
\[
 \|y\|_{2\text{-var};I_j}
 \leq2\|c\|_{2\text{-var};I_j}.
\]
Near-optimal local partitions concatenate into one partition of mesh at
most $\mesh(D)$, so
\[
 \sum_j\|c\|_{2\text{-var};I_j}^2
 \leq w_2(c,\mesh(D))^2.
\]
Combining the three inequalities proves (2.1).
\end{proof}

\begin{hypothesis}[Simple interpolation, BG]\label{hyp:bg}
For every $\varepsilon>0$ and every Jordan parametrization $c$, there is a
partition $D$ with $\mesh(D)<\varepsilon$ such that the parametrized affine
interpolant $c^D$ is a simple polygonal Jordan parametrization.
\status{CITED: Boedihardjo--Geng, Theorem 2.2; exact parameter-mesh and
parametrized-interpolant wording must remain pinned}
\end{hypothesis}

Under Interface~\ref{hyp:bg}, choose $D_n$ with mesh tending to zero and
put $P_n=c^{D_n}$.  The preceding proved lemma gives
\[
 \|P_n-c\|_{2\text{-var}}\to0,
 \qquad P_n\to c\text{ uniformly},                       \tag{2.2}
\]
the uniform statement following because $P_n(0)=c(0)$.

\section{The regular-smoothing interface}

\begin{unknown}[Regular parametrized corner smoothing]\label{unk:smooth}
For every parametrized simple polygon $P\colon\T\to\R^2$ and every
$\varepsilon>0$, construct a $C^\infty$ \emph{regular} Jordan embedding
$Q\colon\T\to\R^2$ such that
\[
 \|Q-P\|_{1\text{-var}}<\varepsilon,                    \tag{3.1}
\]
with compatible collars at every exit point and at the periodic seam.
Here regular means that the chosen parametrization has nonzero velocity
everywhere if that is part of the external smooth-Jordan convention.
\status{UNKNOWN at exact regularity; the geometric rounding is plausible,
but the collar/jet/nonzero-speed parametrization has not been formalized}
\end{unknown}

Assuming Obligation~\ref{unk:smooth}, choose $c_n$ from $P_n$ with errors
$\varepsilon_n\downarrow0$.  Since
$\|x\|_{2\text{-var}}\leq\|x\|_{1\text{-var}}$, (2.2) gives
\[
 \|c_n-c\|_{2\text{-var}}\to0,
 \qquad c_n\to c\text{ uniformly}.                      \tag{3.2}
\]

\section{Uniform local variation}

Define, for $0<\delta\leq\tfrac12$,
\[
 \nu(\delta)=
 \sup_n\sup_{0\leq t-s\leq\delta}
       \|c_n\|_{2\text{-var};[s,t]},                    \tag{4.1}
\]
where intervals are interpreted on the lifted periodic parameter.

\begin{lemma}[Whole-family local modulus]\label{lem:nu}
Under (3.2), $\nu(\delta)\to0$ as $\delta\downarrow0$.
\status{PROVED, conditional only on the approximants supplied above}
\end{lemma}

\begin{proof}
Put $e_n=\|c_n-c\|_{2\text{-var};[0,1]}$.  Restriction, the seam insertion
argument, and Lemma~\ref{lem:local-modulus} give
\[
 \|c_n\|_{2\text{-var};[s,t]}
 \leq\sqrt2\bigl(w_2(c,t-s)+e_n\bigr).                  \tag{4.2}
\]
Given $\eta>0$, first choose $N$ so the right side is small for $n\geq N$,
specifically $\sqrt2e_n<\eta/2$, then choose $\delta$ using $V_0(c)$ so
$\sqrt2w_2(c,\delta)<\eta/2$.  For the finite prefix $n<N$, each
regular $C^1$ map is Lipschitz, say with constant $L_n$, and
\[
 \sum_i|\Delta c_n|^2
 \leq(\max_i|\Delta c_n|)\sum_i|\Delta c_n|
 \leq L_n^2\delta^2.
\]
Taking the finite maximum proves the claim.
\end{proof}

\section{Winding control and primitive compactness}

\begin{proposition}[Embedded-arc winding estimate]\label{prop:arc}
Assume Interface~\ref{hyp:current}.  Let
$a\colon[u,v]\to\R^2$ be an injective rectifiable arc.  Close it by the
oriented chord from $a(v)$ to $a(u)$, obtaining $L$.  Then
\[
 \int_{\R^2}|\Ind(L,z)|\,dz
 \leq\frac\pi4\|a\|_{2\text{-var};[u,v]}^2.             \tag{5.1}
\]
\status{PROVED conditional on the STANDARD-ASSUMED planar interfaces below}
\end{proposition}

\begin{proof}
The components $(\alpha_j,\beta_j)$ on which the arc leaves its closing
chord form disjoint excursions.  Close each excursion by its subchord,
forming a Jordan loop $E_j$ with bounded domain $\Omega_j$.  The parameter
intervals are disjoint, hence a single ordered partition gives
\[
 \sum_j\diam(E_j)^2
 \leq\|a\|_{2\text{-var};[u,v]}^2.                       \tag{5.2}
\]
The planar isodiametric inequality gives
\(
 |\Omega_j|\leq(\pi/4)\diam(E_j)^2.
\)
As integral currents, the original loop is the mass-convergent signed sum
of the excursion loops; the residual cycle supported on the chord line
vanishes because one-forms restricted to a line have primitives.  Planar
filling uniqueness identifies the signed filling multiplicity with
$\Ind(L,\cdot)$.  The triangle inequality and (5.2) prove (5.1).
\end{proof}

\begin{hypothesis}[Excursion and planar-current package]\label{hyp:current}
For an injective rectifiable arc and its closing chord: the off-chord
parameter set is a countable union of open intervals whose endpoints land
on the chord; every excursion plus its subchord is Jordan; its diameter is
bounded by the diameter of the excursion together with its chord endpoints;
and the associated signed boundaries are mass-summable and add to the
original boundary.  Moreover, compactly supported planar integral cycles
have unique compactly supported two-dimensional fillings, whose integer
multiplicity agrees almost everywhere with winding number.  Finally, the
planar isodiametric inequality is available.
\status{STANDARD-ASSUMED; each clause requires a source or formal proof}
\end{hypothesis}

Let
\[
 \lambda_0=\frac{x\,dy-y\,dx}{2},\qquad d\lambda_0=dx\wedge dy,
\]
and define normalized lifted primitives
\[
 F_n(0)=0,
 \qquad F_n(t)-F_n(s)=\int_{[s,t]}c_n^*\lambda_0.         \tag{5.3}
\]

\begin{hypothesis}[Generalized Green formula]\label{hyp:green}
For every closed rectifiable planar curve $L$,
\[
 \int_L\lambda_0=\int_{\R^2}\Ind(L,z)\,dz.              \tag{5.4}
\]
Here $L$ has its displayed orientation, $dx\wedge dy$ is the positive
orientation of the plane, and $\Ind$ uses the corresponding
counterclockwise-positive convention.
\status{CITED: Cuf\'i--Verdera, main theorem at this regularity}
\end{hypothesis}

\begin{proposition}[Primitive compactness]\label{prop:primitive}
Assume Interfaces~\ref{hyp:current} and \ref{hyp:green}.  The functions
$F_n$ have a subsequence
converging locally uniformly on $\R$.
\status{PROVED CONDITIONAL from Lemma~\ref{lem:nu},
Proposition~\ref{prop:arc}, Interfaces~\ref{hyp:current} and
\ref{hyp:green}, and STANDARD-CITED Arzel\`a--Ascoli}
\end{proposition}

\begin{proof}
Uniform convergence places all images in one ball $B(0,R)$.  For
$0\leq t-s\leq\delta$, close $c_n|_{[s,t]}$ by its chord.  Equations
(5.1)--(5.4) and direct chord parametrization yield
\[
 |F_n(t)-F_n(s)|
 \leq \frac\pi4\nu(\delta)^2+\frac R2\nu(\delta).       \tag{5.5}
\]
This is a common modulus tending to zero.  A finite cover of $[0,1]$
bounds $F_n$ uniformly there, so Arzel\`a--Ascoli supplies uniform
convergence on $[0,1]$ along a subsequence.  With $A_n=F_n(1)$,
\[
 F_n(k+r)=kA_n+F_n(r),\qquad k\in\mathbb Z,\quad0\leq r<1,
\]
because the pulled-back integrand is periodic.  This upgrades convergence
to local uniform convergence on $\R$.
\end{proof}

For the alternative Liouville form $\lambda_{\mathrm{AI}}=y\,dx$,
\[
 y\,dx=-\lambda_0+d(xy/2).                               \tag{5.6}
\]
Thus its normalized primitive is
\[
 G_n(t)=-F_n(t)
 +\frac{x_n(t)y_n(t)-x_n(0)y_n(0)}2.                     \tag{5.7}
\]
Equations (3.2), (5.6), and (5.7) prove local uniform convergence of $G_n$
along the same subsequence.
\status{PROVED elementary exact-form conversion}

\section{Conditional rectangle theorem}

\begin{hypothesis}[Asano--Ike approximation criterion]\label{hyp:ai}
The exact primary theorem must be pinned to assert that the following data:
\begin{enumerate}
\item a Jordan parametrization $c$;
\item regular smooth Jordan parametrizations $c_n\to c$ parametrically in
  $C^0$; and
\item locally uniform convergence on $\R$ of normalized lifted primitives
  of the required Liouville form,
\end{enumerate}
imply a prescribed-angle rectangle conclusion for every
$\theta\in(0,\pi)$.
\status{UNKNOWN/CITED-UNVERIFIED interface; a pinned version, theorem/page,
and faithful exact-hypothesis paraphrase must be independently checked}
\end{hypothesis}

\begin{unknown}[Nondegeneracy and boundary order]\label{unk:nd}
Verify from the primary theorem and its definitions that the
$\theta$-rectangle conclusion consists of four pairwise distinct points of
$c(\T)$, in the required cyclic boundary order, and is off the diagonal in
the relevant configuration space.  Also verify that $\theta$ is the angle
between the diagonals and that $\theta=\pi/2$ therefore gives a positive-size
classical square.
\status{UNKNOWN in this formalization; no nondegenerate square may be
asserted from Interface~\ref{hyp:ai} alone}
\end{unknown}

\begin{hypothesis}[Explicit nondegenerate conclusion $\mathrm{AI}_{ND}$]
\label{hyp:aind}
For $\theta=\pi/2$, the conclusion imported from Asano--Ike supplies four
pairwise distinct points of $c(\T)$ in cyclic boundary order forming a
nondegenerate Euclidean rectangle whose diagonals meet at angle $\pi/2$.
\status{UNKNOWN; this is the precise proposition needed from the source}
\end{hypothesis}

\begin{lemma}[Perpendicular-diagonal rectangle]
A nondegenerate Euclidean rectangle whose diagonals are perpendicular is a
positive-size square.
\status{PROVED}
\end{lemma}
\begin{proof}
For adjacent side vectors $u,v$ with $u\perp v$, the diagonal vectors are
$u+v$ and $u-v$.  Their scalar product is $|u|^2-|v|^2$; perpendicularity
therefore gives $|u|=|v|>0$.
\end{proof}

\begin{theorem}[Analytic reduction to the cited rectangle criterion]
Assume Interface~\ref{hyp:bg}, Obligation~\ref{unk:smooth},
Interfaces~\ref{hyp:current} and \ref{hyp:green}, and the unverified
Interface~\ref{hyp:ai}.
If $c$ is a Jordan parametrization with
$V_0(c)$, then $c$ satisfies the approximation and primitive-convergence
hypotheses listed in Interface~\ref{hyp:ai}.  Consequently $c$ has
\emph{exactly the conclusion stated by the verified version} of the cited
Asano--Ike theorem.
\status{PROVED CONDITIONAL REDUCTION; not a nondegenerate rectangle claim}
\end{theorem}

\begin{proof}
Interface~\ref{hyp:bg} and the proved interpolation estimate produce the
simple polygons $P_n$.  Obligation~\ref{unk:smooth} produces regular smooth
Jordan $c_n$ satisfying (3.2).  Lemma~\ref{lem:nu} gives the common local
variation modulus.  Proposition~\ref{prop:primitive}, Interface
~\ref{hyp:green}, and (5.6)--(5.7) give the locally uniformly convergent
lifted primitives.  These are precisely the data listed in Interface
~\ref{hyp:ai}; therefore only the verified conclusion of that interface may
be imported.
\end{proof}

\begin{corollary}[Positive square, conditional on nondegeneracy]
Under all hypotheses of the preceding theorem and assuming the explicit
Interface~\ref{hyp:aind}, $c$ inscribes a positive-size square with four
distinct vertices in cyclic boundary order.
\status{CONDITIONAL; not available while Interface~\ref{hyp:aind} is UNKNOWN}
\end{corollary}

\section{Machine-facing dependency ledger}

\begin{center}
\begin{tabular}{>{\raggedright\arraybackslash}p{0.25\textwidth}
>{\raggedright\arraybackslash}p{0.19\textwidth}
>{\raggedright\arraybackslash}p{0.44\textwidth}}
Component & Status & Exact interface used \\
\hline
Fine-mesh/local modulus & PROVED & Definition $V_0$; seam factor $\sqrt2$ \\
PL interpolation error & PROVED & Every partition; constant $\sqrt8$ \\
Simple PL Jordan choice & CITED & BG theorem, parameter-mesh version \\
Regular smoothing & UNKNOWN & regular $C^1/C^\infty$ embedding; (3.1) \\
Whole-family modulus & PROVED & tail-first then finite-prefix quantifiers \\
Arc winding estimate & CONDITIONAL & excursion/current package \\
Generalized Green & CITED & arbitrary closed rectifiable planar curve \\
Primitive compactness & PROVED CONDITIONAL & current/Green interfaces,
common modulus, Arzel\`a--Ascoli \\
Liouville conversion & PROVED & equations (5.6)--(5.7) \\
AI approximation criterion & UNKNOWN/CITED-UNVERIFIED & version and exact interface must be pinned \\
Four distinct vertices & UNKNOWN & off-diagonal/noncollapse conclusion \\
$\theta=\pi/2$ square & CONDITIONAL & follows only after previous row \\
\end{tabular}
\end{center}

\section{Source identifiers (pinning incomplete)}

The imported interfaces have the following source candidates.  A version
and PDF-page reference is part of a complete pin; any missing page below is
therefore an explicit verification blocker rather than an immutable pin.
\begin{itemize}
\item Boedihardjo--Geng: arXiv:1309.1576v2, Theorem 2.2, PDF pp.~6--7,
  \url{https://arxiv.org/abs/1309.1576v2}.
\item Cuf\'i--Verdera: arXiv:1306.6832v1, main theorem, PDF p.~1
  (specialized to the relevant linear test form),
  \url{https://arxiv.org/abs/1306.6832v1}.
\item Asano--Ike: arXiv:2412.21057v3, Theorem 1.1, PDF p.~2;
  nondegeneracy discussion at PDF pp.~2--3 and the off-diagonal argument in
  the proof of Theorem 4.1 at PDF p.~19 remain subject to independent review,
  \url{https://arxiv.org/abs/2412.21057v3}.
\end{itemize}

\section{Scope}

This draft neither proves the full Square Peg conjecture nor promotes PR
\#122 beyond \texttt{sketch}.  It does not use PRs \#115, \#118, or \#119 as
premises.  There is currently no problem-specific
\texttt{problems/square-peg/RULES.md}; consequently the root repository rules
are the only status protocol available, and this missing problem-level
protocol is itself a governance gap rather than permission to promote.

\end{document}
